\documentclass[11pt]{article}
\usepackage[margin=1in]{geometry}
\usepackage{amsmath,amssymb,amsthm,mathtools}
\usepackage{enumitem}
\usepackage{array,longtable,booktabs}
\usepackage[hidelinks]{hyperref}
\usepackage{microtype}

\newtheorem{theorem}{Theorem}[section]
\newtheorem{lemma}[theorem]{Lemma}
\newtheorem{corollary}[theorem]{Corollary}
\theoremstyle{definition}
\newtheorem{definition}[theorem]{Definition}
\newtheorem{remark}[theorem]{Remark}

\newcommand{\ALCIRCC}{\mathcal{ALCI}_{\mathrm{RCC5}}}
\newcommand{\EQ}{\mathsf{EQ}}
\newcommand{\PP}{\mathsf{PP}}
\newcommand{\PPI}{\mathsf{PPI}}
\newcommand{\PO}{\mathsf{PO}}
\newcommand{\DR}{\mathsf{DR}}
\newcommand{\Rel}{\mathsf{Rel}}
\newcommand{\Comp}{\operatorname{Comp}}
\newcommand{\St}{\operatorname{St}}
\newcommand{\core}{\operatorname{core}}
\newcommand{\U}{\operatorname{U}}
\newcommand{\code}[1]{\texttt{#1}}

\title{Total Transport for the Cluster-Quasimodel Transition System:\\
Round-18, Implementing M1--M6 --- with the Birth Flip}
\author{Claude (Fable~5, Anthropic)\thanks{Round-18 delta on rounds
15--17, implementing the twelfth review's repair package M1--M6
(\code{papers/gpt5.5\_round17\_review/}).  Everything not amended is
inherited.  Ledger going in: twelve reviews --- eleven defects with
witnesses, one no-counterexample verdict.  This delta is unreviewed;
presume it contains the next finding.  Designated weak points:
Section~\ref{sec:invalidate}.}}
\date{July 13, 2026}

\begin{document}
\maketitle

\begin{abstract}
\noindent
The twelfth review rejected round 17 at a precise seam: the ordered-pair
transport rule covered only pairs whose endpoints are both updated by
the ambient transition, omitting mixed pairs with a parent-copy
endpoint, so the pair condition could miss operational pairs.  This
delta implements the review's repair package with one sharpening that
working out the path semantics makes unavoidable: the total endpoint
update $\U$ has \emph{three} clauses, not two.  Besides the review's
ambient clause (recorded steering rows carried forward) and resident
clause (catalogue read-off), cross-branch occurrences require a
\emph{birth-flip} clause: a younger occurrence's entry to an older
separator member is the converse of the steering value that the
\emph{older} side received at the younger's birth.  With $\U$ total,
pairs are transported in lockstep, the inclusion theorem is re-proved by
a five-way endpoint-class analysis, and an order-invariance lemma gives
the least-common-ancestor path semantics the review asked for.  The
review's syntactic demands are also implemented: injective slot maps
(whose necessity the new harness independently re-derived when
non-injective maps produced conflicting pattern writes), the
continued/core consistency invariant, and a sibling-ordering convention
completing the witness map.  The harness architecture answers Finding
17.3 structurally: rule-computed states, derived without ever reading
the realized frame, are cross-checked against frame truth on genuinely
branching unfoldings --- 150 runs, zero divergences, 2{,}820 mixed pairs
and 140 birth-flip entries exercised, a constructed negative control,
and hash-seed-deterministic output.  Open and unchanged: the width
budgets (F6) and uniformization (W2$'$).  Status: \emph{strongly
supported, not certified}.
\end{abstract}

\tableofcontents

\section{Scope, and one sharpening of the work order}\label{sec:scope}

The twelfth review's M1 proposed a total endpoint update with two
clauses: ambient endpoints keep their $T$-images; parent-copy endpoints
take birth-adjacent catalogue states.  Implementing M5 (the path
semantics for ``later gluings elsewhere'') shows a third case is
forced.  Consider an occurrence $y$ born in a copy $A$ attached
\emph{after} some copy $k$ elsewhere in the tree, and a later gluing
$h$ whose separator contains a member $s\in V_k$.  The value
$\rho(y,s)$ was not assigned when $k$ was attached ($y$ did not exist)
nor by $y$'s own steering at any later step: it was assigned \emph{at
$y$'s birth}, when $s$ was a pre-existing occurrence steered at $A$ and
received the value $f_A(\St_A(s))(y)$ toward the fresh $y$.  The entry
$y\to s$ is therefore the \emph{converse of the older side's steering
value at the younger's birth}.  We call this the \emph{birth flip}.  It
is the least-common-ancestor role change that the twelfth review's
Finding 17.2 said was undefined; making it a named clause of $\U$ is
the content of this delta beyond the literal work order.

\section{The total endpoint update}\label{sec:U}

Fix a certificate in the normal form of round 17 (Q3$'$), tightened by
Section~\ref{sec:syntax} below.  For occurrences $x,s$ with $s$ a
member of the separator $\sigma_h$ of a gluing $h$ (so $s$ is a member
of $h$'s parent copy $k$, or a core reference), the entry
$\rho(x,s)$ is given by exactly one of:

\begin{definition}[$\U$: three clauses]\label{def:U}
Let $b(x)$ denote the attachment step of $x$'s birth copy, with core
occurrences born at the base.
\begin{description}[leftmargin=2.2em,style=nextline]
\item[(U-core)] $s\in\core$: the persistent core-row entry of $x$
(in-pattern at $x$'s birth, by (N3)).
\item[(U-res)] $x\in V_k$ (a co-member of $s$'s copy): the pattern
value $N_k(x,s)$, a catalogue value.
\item[(U-down)] $b(x)<b(s)$ and $x\notin V_k$: $x$ was steered when $k$
was attached; the entry is the recorded steering value
$f_{g_k}(\St_{g_k}(x))(s)$.
\item[(U-flip)] $b(x)>b(s)$: at $x$'s birth step, every pre-existing
occurrence is either a member of $x$'s birth copy or steered there
(there is no third possibility --- Lemma~\ref{lem:dichotomy}); the
entry is, respectively, the converse of the in-pattern value
$N_{A}(\,\cdot\,,x)$, or the converse of the recorded steering value
$f_{g_A}(\St_{g_A}(s))(x)$.
\end{description}
$\U_{\!*\to h}(x)$ is the state of $x$ at $h$: its type, persistent core
row, the $\sigma_h$-row assembled from the clauses above, and its
certificate-determined safety signature.
\end{definition}

\begin{lemma}[Birth dichotomy]\label{lem:dichotomy}
At every attachment step, every pre-existing occurrence is a member of
the new copy (equivalently, of its separator, by (N1)) or is steered at
that step.  Hence every pair of occurrences receives its value exactly
once --- in-pattern if ever co-patterned, else by steering at the
younger's birth --- and (U-down)/(U-flip) always find a recorded value.
\end{lemma}

\begin{proof}
The construction steers precisely the pre-existing occurrences outside
the new copy; (N1) says the pre-existing members of the new copy are
exactly the separator.  Single assignment is round-17's no-overwrite
lemma, unchanged.  The recorded value exists because the younger's
birth step steered (or co-patterned) the older endpoint, by the first
sentence.
\end{proof}

\begin{lemma}[$\U$ is total, finite, and catalogue-computable
(M1/M3)]\label{lem:total}
The four clauses are exhaustive and mutually exclusive on the pairs
$(x,s)$ with $s\in\sigma_h$, and each reads finitely many values from
the catalogue, the certificate's steering functions, and states already
computed at earlier steps.  The state alphabet is unchanged from round
17.
\end{lemma}

\begin{proof}
Exhaustiveness: $s\in\core$ or $s\in V_k\setminus\core$; in the latter
case $x\in V_k$ or not, and if not, $b(x)\neq b(s)$ (distinct copies)
splits into (U-down)/(U-flip); Lemma~\ref{lem:dichotomy} supplies the
recorded value in both.  Computability: (U-core), (U-res) are catalogue
lookups; (U-down), (U-flip) are lookups into steering values
$f(\varsigma)(\cdot)$ at states $\varsigma$ computed at strictly
earlier steps --- well-founded in the attachment order.
\end{proof}

\begin{definition}[Pair transport under $\U$ (M2)]\label{def:P3}
Replace round-17's (P3) by: if $(\varsigma,\varsigma',W)$ is reachable
before a step and both endpoints persist, then
$\bigl(\U(\varsigma),\U(\varsigma'),W\bigr)$ is reachable after it,
together with its converse.  Creation (P2) is retained \emph{in both
orientations}: at a birth step, for every steered or co-patterned
pre-existing $x$ and fresh $y$, both $(\St(x),\St(y),V)$ and
$(\St(y),\St(x),V^{-1})$ enter the pair set, where $V$ is the assigned
or in-pattern value.  Co-patterned seeds (P1) are unchanged.
\end{definition}

\section{Order invariance and the path semantics (M5)}\label{sec:order}

\begin{lemma}[Order invariance]\label{lem:order}
The value of every pair, and hence every state $\U_{\!*\to h}(x)$,
depends only on the attachment \emph{tree} (which copies attach where,
with which slot maps) and not on the interleaving order in which
incomparable branches are expanded.
\end{lemma}

\begin{proof}
By Lemma~\ref{lem:dichotomy}, the value of a pair $(x,y)$ is fixed at
the later of the two birth steps, by data local to that step: the
younger's birth copy, slot map, and the older's state there.  The
older's state at that step is, inductively, a function of pair values
fixed at even earlier birth steps along the relevant path.  No step
touches a pair whose endpoints both pre-exist (no overwrite), so
expanding an incomparable branch neither changes existing values nor
the states read at later births.  Induction on the tree.
\end{proof}

\begin{corollary}[LCA path semantics]\label{cor:lca}
For occurrences $y,z$ and a gluing $h$, the pair state
$(\St_h(y),\St_h(z),\rho(y,z))$ is determined by the attachment-tree
paths from the birth copies of $y$ and $z$ through their least common
ancestors to $h$'s parent, with (U-down) applying along descents on the
elder side and (U-flip) applying at and above the younger's birth.
``Later gluing elsewhere in the tree order'' (round-17's undefined
phrase) is thereby replaced: the fixed point closes under one-step
attachment, and Lemma~\ref{lem:order} makes the result path-functional.
\end{corollary}

\section{The inclusion theorem, re-proved (M4)}\label{sec:incl}

\begin{theorem}[Inclusion, total-transport form]\label{thm:incl}
In every finite unfolding prefix of a certificate in normal form: at
every gluing, every existing occurrence's state and every ordered pair
$(\St(y),\St(z),\rho(y,z))$ lie in the fixed point computed with seeds
(P1), two-sided creation (P2), and $\U$-transport
(Definition~\ref{def:P3}).
\end{theorem}

\begin{proof}
Induction over attachment steps, as in round 17, with case (c) replaced
by a five-way endpoint-class analysis.  At a step, classify each
pre-existing endpoint as \emph{resident} (member of the new copy, i.e.\
separator or core) or \emph{steered}; steered endpoints subdivide by
Definition~\ref{def:U} into down and flip histories, but their state
update is uniformly $\U$.  For a pair of pre-existing occurrences, the
five combinations are:
\emph{steered/steered} --- transported by Definition~\ref{def:P3},
both endpoints $\U$-updated, $W$ unchanged;
\emph{steered/resident} and \emph{resident/steered} --- transported
likewise: $\U$ is total, so the resident endpoint's update (U-res) is a
legitimate transport step (this is exactly the mixed case round 17
omitted and Finding 17.1 witnessed);
\emph{resident/resident} --- co-patterned at this step, hence a (P1)
seed, whose value agrees with any pre-existing value by gluing
coherence and no-overwrite.
Pairs with a fresh endpoint are (P2) creations, in both orientations.
Fresh/fresh pairs are (P1) seeds.  States: residents take
catalogue-computable states; steered occurrences take
$\U$-images of inductively-reachable states.  Order invariance
(Lemma~\ref{lem:order}) licenses treating ``the next gluing'' in tree
order without loss of generality.
\end{proof}

\begin{corollary}
Round-16's Theorem 3.2 (deterministic canonical completion) holds with
(S4) checked on the fixed point of Definition~\ref{def:P3}: every
operational ordered pair --- including the mixed ambient/resident and
cross-branch flip pairs --- is present with its current states and
realized mutual value.  The twelfth review's witness geometry is
rejected: the pair $(\St_h(y),\St_h(z),\PP)$ is transported, S4 demands
$f(\St_h(y))(b)\in\Comp(\PP,f(\St_h(z))(b))$, and the bad $(\PO,\DR)$
choice violates $\Comp(\PP,\DR)=\{\DR\}$ (WP32 part B2).
\end{corollary}

\section{Certificate syntax (17.4), the overlap invariant (17.5), and
the sibling convention (R6)}\label{sec:syntax}

\begin{definition}[Port discipline]\label{def:ports}
Every template's ports are partitioned into core ports, inherited
slots, and fresh ports.  Every gluing declaration maps the child's
inherited slots \emph{injectively} into the parent's non-core ports
plus core; fresh occurrence identities are allocated per copy; non-core
identity reuse occurs only through slot maps (connected traces).  (N1)
of round 17 is thereby a syntactic consequence of the declarations, not
a run property.
\end{definition}

\begin{remark}[Injectivity is not decorative]
WP32's own development re-derived its necessity: with a non-injective
slot map, one occurrence occupies two pattern ports, the instantiation
performs conflicting writes, and the rule-computed states diverge from
frame truth --- the harness caught exactly this before injectivity was
enforced, after which agreement is exact ($150$ runs, $0$
divergences).  This is the twelfth review's 17.4, confirmed
empirically.
\end{remark}

\begin{remark}[Continued/core overlap (17.5)]
Core members are both persistent core-row indices and separator
members at every step.  The two sources agree by construction:
(U-core) and (U-res)/(U-down)/(U-flip) coincide on core columns because
core values are in-pattern at every copy (N3) and never overwritten.
We state this as an explicit invariant so a mechanization can discharge
it once.
\end{remark}

\begin{remark}[Sibling ordering (R6 completion)]
Children of a copy are ordered lexicographically by (source member
index, demand index in closure order); demands designating the same
gluing edge at the same source share one copy.  With this convention
and the witness map of round 17, the unfolding is a function of the
certificate, including its tree structure.
\end{remark}

\section{Machine checks (WP32)}\label{sec:wp}

\code{verification/python/wp32\_round18\_total\_transport.py}, all
PASS, with the harness architecture the twelfth review's Finding 17.3
demanded: rule-computed states (never reading the frame) cross-checked
against frame truth on genuinely \emph{branching} unfoldings.
\textbf{A} (M1/M5): 150 random branching unfoldings, \emph{zero}
$\U$-clause divergences between rule and frame.
\textbf{B} (M2/17.1 regression): 2{,}820 mixed ambient/resident and
flip pairs exercised, 140 birth-flip entries; the WP31 geometry's mixed
pair is transported and S4 rejects the bad row.
\textbf{C} (M6): a constructed adversarial steering violating S4
produces closure failures in $17/41$ runs --- the negative control
fires (round-17's harness never sampled one).
\textbf{D} (17.3): no set-iteration choice points; a result digest is
printed and is stable across \code{PYTHONHASHSEED} values.

\section{Response map}\label{sec:resp}

\begin{center}
\begin{tabular}{p{0.12\linewidth}p{0.80\linewidth}}
\toprule
M1 & $\U$ total on all pre-existing endpoints
(Definition~\ref{def:U}), with the birth-flip clause added beyond the
work order's two cases (Section~\ref{sec:scope}).\\
M2 & Pair transport under $\U$ with converses; two-sided creation
(Definition~\ref{def:P3}).\\
M3 & Lemma~\ref{lem:total}.\\
M4 & Theorem~\ref{thm:incl}: five-way endpoint-class analysis replacing
``both endpoints update by $T$''.\\
M5 & Lemma~\ref{lem:order} and Corollary~\ref{cor:lca}: one-step
closure plus order invariance; the LCA characterization derived, the
undefined phrase eliminated.\\
M6 & WP32: branching trees, mixed transport exercised and counted,
constructed negative S4 control, deterministic output.\\
17.4 & Definition~\ref{def:ports} (with empirical confirmation).\\
17.5 & The overlap invariant remark.\\
R6 & The sibling convention remark.\\
\bottomrule
\end{tabular}
\end{center}

\section{What would invalidate this delta}\label{sec:invalidate}

\begin{enumerate}[label=(Z\arabic*),leftmargin=2.6em]
\item \textbf{A gap in the birth dichotomy or the flip clause}: a
relative position of $(x,s)$ not covered by
Definition~\ref{def:U} --- e.g.\ if some construction step neither
co-patterns nor steers a pre-existing occurrence (this would contradict
Lemma~\ref{lem:dichotomy}'s first sentence, which is definitional here,
but an extraction-side presentation might demand such a step).
\item \textbf{Order invariance under declarations richer than the
normal form}: Lemma~\ref{lem:order} leans on no-overwrite and
injective slot maps; exotic-but-legal certificates are the place to
look.
\item \textbf{The standing open items}: F6 (width budgets), W2$'$
(uniformization), and the single external input (RCC5 patchwork).
\end{enumerate}

\section{Status}

Twelve reviews: eleven defects with witnesses, one no-counterexample
verdict.  This delta implements the twelfth review's full repair
package, sharpened by the birth-flip clause that the path semantics
forces, and its harness now checks the \emph{text} against frame truth
by construction --- the mismatch class that produced Findings 17.1 and
17.3 is structurally excluded from the toolchain, not just from this
round.  The delta is unreviewed; presume the next finding.  The
soundness half remains, as of this writing, fully written and
machine-cross-checked; the completeness half retains its two named open
lemmas.  The decidability theorem for $\ALCIRCC$ remains
\textbf{strongly supported, not certified}.

\begin{thebibliography}{9}
\bibitem{review12} GPT-5.5 (OpenAI).  \emph{Formal Response on the
Round-17 Transition-System Delta}, July 13, 2026.
\code{papers/gpt5.5\_round17\_review/}.

\bibitem{round17} Claude (Fable 5).  \emph{The Explicit Transition
System for Certified Joint Steering}, July 13, 2026.
\code{papers/fable5\_round17/}.

\bibitem{round16} Claude (Fable 5).  \emph{Certified Joint Steering for
Cluster Quasimodels}, July 13, 2026.  \code{papers/fable5\_round16/}.

\bibitem{renznebel} J.~Renz and B.~Nebel.  On the complexity of
qualitative spatial reasoning.  \emph{Artificial Intelligence} 108
(1999).
\end{thebibliography}

\end{document}
